\printedpage{4}
% head: A. Grothendieck : Dérivateurs, Ch. V, Édition des Universités de Montpellier II et Paris VII
% foot: Version 9 septembre 2001

(This is not exciting either \dots)

Let
\[
\varepsilon=\{-1,+1\}
\]
be the ordered set, \textit{discrete,} reduced to the two elements $-1$, $+1$, and let us define inclusion functors (fully faithful)
\[
\begin{aligned}
\alpha_n:\varepsilon\longrightarrow\Phi_n\quad & \alpha_n(-1)=-n\\
& \alpha_n(1)=n,
\end{aligned}
\]
these functors are compatible with the $\varphi_{n,m}$, whence % [sic] compatible -> compatibles, accord avec « ces foncteurs »
\[
\alpha_\infty:\varepsilon\longrightarrow\boldsymbol{\Phi}\in\mathrm{ProCat},
\]
whence for every $X$ a functor
\[
\kappa:\uHom(\boldsymbol{\Phi},X)\longrightarrow\uHom(\varepsilon,X)\simeq X\times X
\]
i.e.

\msp{page 5}
\[
\begin{aligned}
\uCh(X) &\longrightarrow X\times X\\
c &\longrightarrow (\sigma(x),\beta(x)).
\end{aligned}
\]

\texttt{[rather} $(\sigma(c),\beta(c))$\texttt{]} For every path $c$ in $X$ i.e. a diagram $(*)$ (satisfying the condition of degeneracy at infinity), we have a source or origin of the path, and a target or an \textit{endpoint}
\[
\left\{
\begin{aligned}
\sigma(x) &= A_i\quad\text{for } i\leq -\mathrm{long}(x)\\
\beta(x) &= A_i\quad\text{for } i\geq \mathrm{long}(x)
\end{aligned}
\right.
\]
\texttt{[replace} $x$ \texttt{by} $c$\texttt{; for long}$(c)$ \texttt{suitably defined]}. The fiber of the preceding functor at $x$, $y$ will be denoted
\[
\uCh(X;x,y)\qquad\text{or simply }\uCh(x,y),
\]
\textit{category of paths from $x$} to \textit{$y$.}

\printedpage{5}
% head: A. Grothendieck : Dérivateurs, Ch. V, Édition des Universités de Montpellier II et Paris VII
% foot: Version 9 septembre 2001

\section*{2\quad Calculation of a category of fractions $X\Sigma^{-1}$}

Before continuing the formalism of paths (composition etc.), I am going to test the usefulness of the notion of category of paths, for the ``calculation'' of the $X\Sigma^{-1}$, where
\[
\Sigma\subset\mathrm{Fl}(X),\qquad\Sigma\text{ contains the identities}
\]
is a subset, otherwise \textit{arbitrary}, of $\mathrm{Fl}(X)$. We know that $\mathrm{Ob}\,X\Sigma^{-1}=\mathrm{Ob}\,X$, and everything comes down to defining the
\[
\Hom_{X\Sigma^{-1}}(x,y),\qquad x,y\in\mathrm{Ob}\,X,
\]
\msp{page 6}
and for $x,y,z\in\mathrm{Ob}\,X$ the composition
\[
\Hom_\Sigma(x,y)\times\Hom_\Sigma(y,z)\xrightarrow{(u,v)\mapsto vu}\Hom_\Sigma(x,z),
\]
(where we write $\Hom_\Sigma$ instead of $\Hom_{X\Sigma^{-1}}$, to lighten the notation).

Using $\Sigma$, we define a full subcategory
\[
\uCh_\Sigma(X)\hookrightarrow\uCh(X),
\]
(\footnote[3]{The set of its objects is denoted $\mathrm{Ch}(X)$. Its elements are the $\Sigma$-paths from $x$ to $y$ \texttt{[from} $x$ \texttt{to} $y$\texttt{?]}.}) by the condition on $c\in\mathrm{Ch}(X)$: For $i$ odd, \texttt{[rather even]} i.e. $i-1\xleftarrow{u_i}i$ ($u_i$ is called ``retrograde''), $c(u_i)\in\Sigma$, i.e. the retrograde ``segments'' in the path $c$ must be in $\Sigma$. We similarly define the full subcategories
\[
\uCh_\Sigma(X;x,y)=\{c\in\uCh(X;x,y)\mid c\in\uCh_\Sigma(X)\},
\]
or again the fiber at $(x,y)$ of the source-target functor for $\uCh_\Sigma(X)$
\[
\uCh_\Sigma(X)\longrightarrow X\times X.
\]

We define a canonical map
\[
\mathrm{Ch}_\Sigma(X,x,y)\longrightarrow\Hom_\Sigma(x,y),
\]
\msp{page 7}
by associating to every path from $x$ to $y$ the $\Sigma$-homomorphism from $x$ to $y$ obtained by \textit{inverting} (in $X\Sigma^{-1}$) the retrograde segments (which is possible, since they are in $\Sigma$), so that we find in $X\Sigma^{-1}$ a ``direct path'' % [sic] `a -> à, accent imprimé comme guillemet ouvrant
\[
x=A_{-n}\longrightarrow\cdots\longrightarrow A_{-1}\longrightarrow A_0\longrightarrow A_1\longrightarrow A_2\cdots A_n=A_{n+1}=\cdots=y.
\]
and we compose all these morphisms.

In fact, we find at once that we thus define a \textit{functor}
\[
\uCh_\Sigma(X)\longrightarrow\uFl(X\Sigma^{-1}),
\]
\texttt{[}$\uFl$ \texttt{category of arrows]} from the fact that if we have a diagram

\printedpage{6}
% head: A. Grothendieck : Dérivateurs, Ch. V, Édition des Universités de Montpellier II et Paris VII
% foot: Version 9 septembre 2001
\[
\begin{tikzcd}
A_{i-1} \arrow[d, "f_{i-1}"] & A_i \arrow[l, "u_i"'] \arrow[d, "f_i"] \\
A'_{i-1} & A'_i \arrow[l, "u'_i"']
\end{tikzcd}
\qquad\qquad (i\text{ even})
\]
with $u_i$, $u'_i$ in $\Sigma$, this diagram, as a diagram in $X\Sigma^{-1}$, is commutative with horizontal arrows \textit{isomorphisms,} and therefore remains commutative on inverting the said arrows. Whence an induced functor
\[
\uCh_\Sigma(X;x,y)\longrightarrow\uFl(\underbrace{X\Sigma^{-1}}_{X'};x,y).
\]

But for any category $X'$, and objects $x$, $y$ in $X'$, the category $\uFl(X',x,y)$ is
\msp{page 8}
a category, \textit{discrete,} having as its set of objects $\Hom_{X'}(x,y)$. Thus the preceding functor in fact defines a canonical map % [sic] précédant -> précédent, adjectif
\[
\boxed{\pi_0\uCh_\Sigma(X;x,y)\longrightarrow\Hom_\Sigma(x,y).}
\]
I say that this map is always bijective. To see this, we are going to \textit{construct} ``from scratch'' the category $X\Sigma^{-1}$, taking as objects those of $X$, and \textit{defining} % [sic] constuire -> construire, lettre omise
\[
\boxed{\Hom_\Sigma(x,y)=\pi_0\uCh_\Sigma(X;x,y).}
\]
We must therefore define the composition
\[
\Hom_\Sigma(x,y)\times\Hom_\Sigma(y,z)\longrightarrow\Hom_\Sigma(x,z).
\]
For this, we are going to define functors ``composition of paths''
\[
\begin{aligned}
\uCh_m(X;x,y)\times\uCh_n(X;y,z) &\longrightarrow \uCh_{m+n}(X;x,z)\\
(c,c') &\longmapsto c'\circ c,
\end{aligned}
\]
but restricting ourselves for simplicity to the case where $m$, $n$ have \textit{the same parity,} so that by ``putting end to end'' $c$ and $c'$ (in this order), we do indeed find a path of type $n+m$.

\msp{page 9}
This operation is ``transposed'' from the pair of functors
\[
\begin{aligned}
\Phi_n &\xrightarrow{\alpha} \Phi_{n+m}\quad i\longmapsto (i-m)\\
\Phi_m &\xrightarrow{\beta} \Phi_{n+m}\quad i\longmapsto (i+n),
\end{aligned}
\]
\printedpage{7}
% head: A. Grothendieck : Dérivateurs, Ch. V, Édition des Universités de Montpellier II et Paris VII
% foot: Version 9 septembre 2001
where we will have
\[
\alpha(\{n\})=\beta(\{-m\})=n-m,
\]
and we will have
\[
\Phi_{n+m}\simeq\Phi_n\sqcup_e\Phi_m,
\]
i.e. $\Phi_{n+m}$ is the ``amalgamated sum'' in (Cat), obtained by ``gluing'' $\Phi_n$ and $\Phi_m$, by identification of the ``target element'' of $\Phi_n$ with the ``source element'' of $\Phi_m$. The composition of paths functor induces \texttt{[in what precedes, from the beginning of page 9, one must interchange} $m$ \texttt{and} $n$\texttt{]}
\[
\uCh_{m,\Sigma}(x,y)\times\uCh_{n,\Sigma}(y,z)\longrightarrow\uCh_{m+n,\Sigma}(x,z),
\]
whence on passing to the $\pi_0$
\[
\Hom_{\Sigma,m}(x,y)\times\Hom_{\Sigma,n}(y,z)\longrightarrow\Hom_{\Sigma,m+n}(x,z),
\]
where we denote, for $a,b\in\mathrm{Ob}\,X$, by
\[
\Hom_{\Sigma,n}(a,b)\overset{\text{def}}{=}\pi_0(\uCh_{m,\Sigma}(a,b)).
\]

\texttt{[rather} $\uCh_{n,\Sigma}$\texttt{]}

\msp{page 10}
To deduce from this a map on the $\varinjlim$
\[
\underbrace{\Hom_\Sigma(x,y)}_{=\ \varinjlim\ \Hom_{\Sigma,m}(x,y)}\quad\times\quad\underbrace{\Hom_\Sigma(y,z)}_{=\ \varinjlim\ \Hom_{\Sigma,n}(y,z)}\quad\longrightarrow\quad\underbrace{\Hom_\Sigma(x,z)}_{=\ \varinjlim\ \Hom_{\Sigma,n+m}(x,z)}
\]
we must prove, for $n\leq n'$, $m\leq m'$, commutativity in
\[
\begin{tikzcd}
\Hom_{\Sigma,m}(x,y)\times\Hom_{\Sigma,n}(y,z) \arrow[r] \arrow[d] & \Hom_{\Sigma,m+n}(x,y) \arrow[d] \\
\Hom_{\Sigma,m'}(x,y)\times\Hom_{\Sigma,n'}(y,z) \arrow[r] & \Hom_{\Sigma,m'+n'}(x,y)
\end{tikzcd}
\]
\texttt{[rather} $\Hom_{\Sigma,m+n}(x,z)$\texttt{,} $\Hom_{\Sigma,m'+n'}(x,z)$\texttt{]} In other words, that for $c:x\longrightarrow y$ and $c':y\longrightarrow z$ two composable $\Sigma$-paths of length $m$ and $n$ respectively, the composite $c'c:x\longrightarrow z$ does not change, ``up to homotopy'' in $\uCh_\Sigma(x,z)$, when we replace $c$, $c'$ by paths $\bar{c}$, $\bar{c}'$ of length $m'$, $n'$ deduced from $c$ and $c'$ by ``degeneracy at the ends''.

More precisely, % [sic] précisement -> précisément, accent omis
\[
\bar{c}'\bar{c}\sim\overline{c'c},
\]
the composite path is homotopic to the one deduced from $c'c$ by ``degeneracy at the ends''.

I have just looked at $n=m=1$, $n'=m'=2$, and found that there is no natural arrow between $\bar{c}'\bar{c}$ and $\overline{c'c}$ - the homotopy sought must
\printedpage{8}
% head: A. Grothendieck : Dérivateurs, Ch. V, Édition des Universités de Montpellier II et Paris VII
% foot: Version 9 septembre 2001
\msp{page 11}
be defined with the help of a ``path'' in $\uCh_{\Sigma,m',n'}(X;x,z)$. In the general case, we are in fact going to find a homotopy
\[
h=h_{(n,m),(n',m')}
\]
``making commutative'' the diagram of functors
\[
\begin{tikzcd}[column sep=large]
\mathcal{A}=\uCh_{\Sigma,m}(X;x,y)\times\uCh_{\Sigma,n}(X;y,z) \arrow[r, "*_{m,n}"] \arrow[d, "i_{m,m'}\times j_{n,n'}"] & \uCh_{\Sigma,m+n}(X;x,z) \arrow[d, "k_{m+n,m'+n'}", ""'{name=K}] \\
\uCh_{\Sigma,m'}(X;x,y)\times\uCh_{\Sigma,n'}(X;y,z) \arrow[r, "*_{m',n'}", ""{name=S}] & \uCh_{\Sigma,m'+n'}(X;x,z)=\mathcal{B},
\arrow[from=K, to=S, bend right, "h"']
\end{tikzcd}
\]
where the $i_{\alpha,\beta}$, $j_{\alpha,\beta}$, $k_{\alpha,\beta}$ are the transition functors
\[
\uCh_{\Sigma,\alpha}(X;a,b)\longrightarrow\uCh_{\Sigma,\beta}(X;a,b)
\]
$(\alpha\leq\beta)$ for the $(x,y)$, $(y,z)$ and $(x,z)$ (in place of $(a,b)$) respectively. The homotopy $h$ will be a path in the category $\uHom(\mathcal{A},\mathcal{B})$, going from
\[
k_{m+n,m'+n'}\circ{*}_{m+n}\longrightarrow{*}_{m'+n'}\circ(i_{m,m'}\times j_{n,n'}).
\]
The existence of such a homotopy means neither more nor less than that the square diagram above is commutative, as a diagram in the category (Cat\,hot) of ``categories up to homotopy,'' whose objects are
\msp{page 12}
the (small) categories, and where we define
\[
\Hom_{\mathrm{Cat\,hot}}(\mathcal{A},\mathcal{B})\overset{\text{def}}{=}\underbrace{\pi_0(\uHom(\mathcal{A},\mathcal{B}))}_{\begin{array}{@{}l@{}}\text{homotopy classes}\\ \text{of functors from }\mathcal{A}\\ \text{to }\mathcal{B}\end{array}}\,.
\]
It is this conclusion that will be useful to us, and not the explicit description of the path $h$. This existence granted (cf. the proof below), we see therefore that if we regard, for $a,b\in\mathrm{Ob}\,X$, the system of categories $\uCh_{\Sigma,r}(X;a,b)$ $(r\in\mathbf{N})$ as an \textit{ind-object of} Cat\,hot\textit{,} which we denote by
\[
\uCh^{\mathrm{ind}}_{\Sigma}\,(X;a,b)\in\mathrm{Ob}\,\mathrm{Ind}(\mathrm{Cat\,hot}),
\]
we will have a canonical morphism in Ind(Cat\,hot)
\[
\uCh^{\mathrm{ind}}_{\Sigma}\,(X;x,y)\times\uCh^{\mathrm{ind}}_{\Sigma}\,(X;y,z)\xrightarrow{\ *\ }\uCh^{\mathrm{ind}}_{\Sigma}\,(X;x,z),
\]
called \textit{``ind-morphism of composition of paths.''} We would have liked to deduce it from an honest \textit{functor} ``composition of paths''
\[
\begin{aligned}
\uCh_\Sigma(X;x,y)\times\uCh_\Sigma(X;y,z) &\xrightarrow{\ ?\ } \uCh_\Sigma(X;x,z),\\
\text{which we would write }(c,c') &\longmapsto c'\circ c\text{ or }c'\cdot c,
\end{aligned}
\]
but I doubt that such a functor exists,